\section{Examples and remarks}\label{sec:examples}

\begin{example}[Determinantal hypersurfaces]\label{ex:determinantal}
Let $d\ge2$, let $l_1,l_2$ be independent linear forms on $\PP^4$ and
$f_1,f_2$ general forms of degree $d-1$, and put
\[
 X_0=\{F_0=0\},\quad F_0=l_1f_2-l_2f_1,\qquad
 \cX=\{(t,x)\in\DD\times\PP^4:\ F_0(x)+tP(x)=0\},
\]
where $P$ is a form of degree $d$ that does not vanish at the points of
$\{l_1=l_2=f_1=f_2=0\}$, and $\DD$ is a small disc. Let $\omega$ be the
restriction to $\cX$ of $\tfrac i2dt\wedge d\bar t+\omega_{FS}$, where
$\omega_{FS}$ is the Fubini--Study form with $[\omega_{FS}]=c_1(\OO(1))$;
thus $g_t$ is the Fubini--Study metric on the hypersurface $X_t$.

The singular set of $X_0$ consists of the $(d-1)^2$ points of
$\{f_1=f_2=0\}$ in the plane $\{l_1=l_2=0\}$, which are ordinary double
points: there $(l_1,l_2,f_1,f_2)$ are local coordinates on $\PP^4$. The
total space $\cX$ is smooth at these points because
$\partial_t(F_0+tP)=P\ne0$, and smooth elsewhere near $X_0$. A small
resolution is the complete intersection
\[
 Y=\{\zeta_0l_1+\zeta_1l_2=0,\ \zeta_0f_1+\zeta_1f_2=0\}
 \subset\PP^4\times\PP^1 ,
\]
of the ample classes $H+L$ and $(d-1)H+L$, where $H$ and $L$ are the
hyperplane classes of the factors. The first divisor is the blow-up of
$\PP^4$ along the plane, on which the second linear system is base-point
free, so $Y$ is smooth for general $f_1,f_2$ by Bertini's theorem, and
$X_0$ is smooth away from the $(d-1)^2$ points because $Y\to X_0$ is an
isomorphism there. By the Lefschetz theorem, applied
twice, $b_1(Y)=0$ and $H^2(Y;\ZZ)=\ZZ H\oplus\ZZ L$, and likewise
$b_2(X_t)=1$. Thus the hypotheses of Theorem~\ref{thm:main} hold, with
$K=\RR H$ and $m=1$. From the class of $Y$,
\[
 \int_YH^3=d,\qquad\int_YH^2L=d-1,\qquad\int_YHL^2=0,
\]
and each exceptional curve is a fibre of the projection to $\PP^4$, so
$H\cdot C_i=0$ and $L\cdot C_i=1$. By \eqref{eq:rank-one} with $D=L$,
\[
 \mu_1=\frac{(d-1)^2}{(d-1)^2/d-0}=d,\qquad
 \lambda^+_{2,1}(t)=\frac{2\pi d}{\log(1/|t|)}\,(1+o(1)).
\]
For $d=5$ this is the split quintic, and the coefficient $10\pi$ is the
one found in \cite{wuebben-conifold} for the Ricci-flat metrics in the
class $c_1(\OO(1))$.
\end{example}

\begin{remark}[Dependence on the metric]\label{rem:metric}
The limits in Theorem~\ref{thm:main} depend on $\omega$ only through the
class $H$. For the Ricci-flat metrics of a polarized Calabi--Yau conifold
transition the same limits are proved in \cite{wuebben-conifold}, with the
logarithm taken along a ray. There the tangent cone at a node is the
Ricci-flat conifold, with radial function homogeneous of degree $2/3$ in
the coordinates of \eqref{eq:normal-form}; the constant
\eqref{eq:link-constant} is $2\pi/3$ and the scale is $|t|^{1/3}$, in
place of $\pi$ and $|t|^{1/2}$, and the product
$c/\log(1/\epsilon)=2\pi/\log(1/|t|)$ is the same. In both cases
$\|\vartheta\|^2_{L^2(L)}=2\pi\alpha\cdot(-\int_{Z_\sigma}\vartheta_Z^2)
=\pi\alpha$ and $\epsilon=|t|^{\alpha/2}$, where $\alpha$ is the degree of
homogeneity of the radial function in the coordinates of
\eqref{eq:normal-form} ($\alpha=1$ here and $\alpha=2/3$ for the
Ricci-flat cone).
\end{remark}

\begin{remark}[The form $G$]\label{rem:G}
The form $G$ of \eqref{eq:gram} agrees with the Hodge--Riemann form of
the class $H$ on the primitive classes and is positive on $\RR H$: it is $-\int_YD^2H$ on classes
with $\int_YDH^2=0$, and $\tfrac12\int_YH^3$ on $H$. The resolution $Y$ enters
only through its cohomology ring; the form $B$ is unchanged by flops of
the curves $C_i$, since $H\cdot C_i=0$, and $\per$ changes by signs, so $\cE$, $\cM$ and the
$\mu_j$ are unchanged.
\end{remark}

\begin{remark}[Scope]\label{rem:scope}
The theorem concerns one-parameter families with smooth total space, so
that every node is smoothed to first order; families in which the nodes
are smoothed at different rates, and the next term in
\eqref{eq:main}, are not addressed. The small two-form eigenvalues are
accompanied by eigenvalues of the same size on exact three-forms and on
their Hodge duals, and we do not determine the complete small spectrum in
degree three.
\end{remark}
